\section{Proofs}\label{app:proofs}
\paragraph{Invariance law.}
Under $B\mapsto BM$, $[A'_i,A'_j]=\sum_{k,l}M_{ki}M_{lj}[A_k,A_l]$, so the commutator matrix becomes
$C(M\otimes M)$, while $P_0(BM)=P_0(B)$. Hence $R\mapsto R(M\otimes M)$, and the bounds follow from
$\sigma(M\otimes M)=\{\sigma_i\sigma_j\}$. For $\varepsilon>0$,
$P_\varepsilon(BM)=BM(M^\top GM+\varepsilon I)^{-1}M^\top B^\top=B(G+\varepsilon M^{-\top}M^{-1})^{-1}B^\top$.

\paragraph{Proposition~\ref{prop:basic}.}
\emph{(a)} The commutator is bilinear and antisymmetric and satisfies Jacobi, so closure makes $\mathfrak a$ a Lie
subalgebra; the Lie subgroup theorem \citep[Ch.~5]{hall2015lie} gives the integrating subgroup. It need not be
closed: an irrational line in a torus is an example.

\emph{(b)} Write $X=\sum x_iA_i$, $Y=\sum y_jA_j$ and $r_{ij}=N[A_i,A_j]$. Then
$\|N[X,Y]\|\le\|x\|\|y\|(\sum\|r_{ij}\|^2)^{1/2}=\|x\|\|y\|K\sqrt{\Lc}$ by Cauchy--Schwarz, and
$\|X\|^2=x^\top Gx\ge\lambda_{\min}\|x\|^2$.

\emph{(c)} Dynkin's form \citep{bonfiglioli2012topics} is
$\log(e^Xe^Y)=\sum_{m\ge1}\frac{(-1)^{m-1}}{m}\sum\frac{[w]}{|w|\prod r_i!s_i!}$, summed over words
$w=X^{r_1}Y^{s_1}\cdots X^{r_m}Y^{s_m}$ with $r_i+s_i\ge1$, where $[w]$ is right-nested.
\begin{itemize}
\item By submultiplicativity, $\|[w]\|\le2^{|w|-1}r^{|w|}$.
\item Let $e_n$ bound $\|N[w]\|$ for $|w|=n$, so $e_1=0$ and $e_2\le\delta r^2$. Write $[w]=[W,Z]$ with $W$ a letter.
\item $[W,P_0Z]$ has both arguments in $\mathfrak a$, so its escape is at most $\delta r\|Z\|\le\delta2^{n-2}r^n$;
and $\|[W,NZ]\|\le2re_{n-1}$. By induction $e_n\le(n-1)2^{n-2}\delta r^n$.
\item Each word therefore contributes at most $\frac{\delta}{4}(2r)^{|w|}/\prod r_i!s_i!$. The full sum is
$\sum_m\frac1m(e^{4r}-1)^m=-\ln(2-e^{4r})$ for $r<\ln2/4$.
\item The length-1 words lie in $\mathfrak a$ and account for $4r$; subtracting them gives the bound.
\item The length-2 words sum to $\frac12[X,Y]$, whose escape is at most $\frac12\delta r^2$; replacing their crude
contribution $4\delta r^2$ gives $\frac12\delta r^2+O(\delta r^3)$.
\item The series converges absolutely on this ball, and the principal logarithm is analytic there
($\|e^Xe^Y-I\|<1$). Both sides agree near $0$, hence on the whole connected ball.
\end{itemize}

\emph{(d)} With $d_t=\hat z_t-z_t$: $d_{t+1}=R_td_t-e_t$ and $\|R_td_t\|=\|d_t\|$.
$\square$

\paragraph{Proposition~\ref{prop:classify}.}
Subalgebras of $\so(d)$ carry the ad-invariant negative-definite form $\operatorname{tr}(XY)$, hence are compact and
reductive.

\emph{(a)} A 2-dimensional reductive algebra is abelian. It is maximal abelian because
$\operatorname{rank}\so(5)=2$, and maximal abelian subalgebras of a compact Lie algebra are conjugate
\citep{adams1969lectures}.

\emph{(b)} A 3-dimensional reductive $\mathfrak h$ is its center plus a semisimple part of dimension $0$ or $3$.
It cannot be abelian (the rank is 2), so $\mathfrak h\cong\su(2)$. Its projections onto the ideals $\su(2)_{L},\su(2)_{R}$ are zero or injective, because
$\mathfrak h$ is simple.
\begin{itemize}
\item If one projection is zero, $\mathfrak h$ is a factor. The factors are ideals, hence invariant under all of
$SO(4)$, and a reflection exchanges them.
\item Otherwise $\mathfrak h$ is the graph of an isomorphism $\su(2)_L\to\su(2)_R$. Any two such graphs differ
by an automorphism of $\su(2)$, which is inner, and $SU(2)\times SU(2)\to SO(4)$ is onto, so all such $\mathfrak h$
are $SO(4)$-conjugate. The stabilizer of $e_4$, $\so(3)$ on $\mathbb R^3\oplus\mathbb R$, is one of them (it is not
an ideal).
\item $SO(4)$ is connected, so it cannot exchange the two ideals; hence three classes under $SO(4)$ and, since a
reflection fixing $e_4$ exchanges them and preserves the diagonal, two under $O(4)$.
\end{itemize}
The Pfaffian is $SO(4)$-invariant and odd under reflections. It vanishes on rank-2 rotations and equals
$\pm\frac12$ on unit elements of $\su(2)_{L/R}$.

\emph{(c)} A 4-dimensional reductive subalgebra is $\su(2)\oplus\mathfrak u(1)$ with the $\mathfrak u(1)$ central.
The centralizer of the diagonal $\so(3)$ is trivial, because $\so(4)\cong\Lambda^2\mathbb R^4$ decomposes as
$\mathbf 3\oplus\mathbf 3$ under it, with no trivial summand; the same holds for its conjugates. So the $\su(2)$ is
a factor, which equals the derived algebra $[\mathfrak h,\mathfrak h]$. A diagonal $\so(3)\subset\mathfrak h$ would
satisfy $\so(3)=[\so(3),\so(3)]\subset[\mathfrak h,\mathfrak h]$, a factor, which is impossible.

\emph{Scope.} The classification is for exactly closed spans. Learned spans are only approximately closed; we assign
them to classes using continuous invariants (Pfaffian, chirality) and the closure-flow check, without a stability
theorem relating approximate closure to distance from a closed class. $\square$

\paragraph{Proposition~\ref{prop:endpoint}.}
Here $g\in SO(3)$ acts on $z_{1:3}$ and fixes $z_4$, so transitions preserve $\|z\|$.

\emph{(a)} Identify $\mathbb R^4$ with the quaternions $\mathbb H$. The exponential of $\su(2)_L$ is left
multiplication by unit quaternions, which acts simply transitively on each sphere $\|z\|=\rho>0$. For
$\|z\|=\|z'\|$, $u=z'z^{-1}$ is the unique unit quaternion with $L_uz=z'$.
\begin{itemize}
\item $q=\log u$ is smooth wherever $u\neq-1$, and $\exp(\log u)=u$ for every branch.
\item $u_{02}=z_2z_0^{-1}=(z_2z_1^{-1})(z_1z_0^{-1})=u_{12}u_{01}$, so composition is exact.
\item Closure is exact, the basis is orthonormal so the Gram penalty is $0$, and $\mathrm{Cov}(z)=I$ in population
(states are $\mathcal N(0,I)$ and rotations preserve that law).
\end{itemize}
The coefficient of a rotation by angle $\theta=\angle(z,z')$ has entries up to $\sqrt2\theta$; only pairs beyond the
bound $0.65$ break exactness.

\emph{(b)} Fix the sphere $S$ and suppose a continuous rule gives $R(z,z')=\rho(q(z,z'))\in SO(3)$ with
$R(z,z')z=z'$ and $R(z_0,z_2)=R(z_1,z_2)R(z_0,z_1)$ for all $z_i\in S$ with pairwise angles $<\epsilon$.
\begin{itemize}
\item The triple $(z,z,z)$ gives $R(z,z)=R(z,z)^2$, so $R(z,z)=I$.
\item For a chain $x_0,\dots,x_n$ with consecutive angles $<\epsilon/3$, let $P=R(x_{n-1},x_n)\cdots R(x_0,x_1)$.
Inserting or deleting a point inside an $\epsilon$-small triangle leaves $P$ unchanged. Two chains with the same
endpoints and pointwise angles $<\epsilon/3$ are joined by such moves (replace one point at a time; each move is a
triangle of diameter $<\epsilon$), so $P$ is unchanged. Fine subdivisions of a continuous path therefore define one
$P$, and by uniform continuity it is constant along a homotopy of paths with fixed endpoints.
\item $S\cong S^2$ is simply connected, so fixing $e\in S$, $s(z)=P(e\to z)$ is well defined, with $s(z)e=z$.
It is continuous, because $s(z')=R(z,z')s(z)$ for nearby $z'$ and $R$ is continuous with $R(z,z)=I$.
\item For a unit $f\perp e$ (in the $\mathbb R^3$ part), $z\mapsto s(z)f$ is orthogonal to $s(z)e=z$: a
nowhere-vanishing tangent field on $S^2$, contradicting the hairy-ball theorem.
\end{itemize}
The rule takes values in the true (diagonal) group; the argument uses no property of $q$ beyond continuity.

\emph{(c)} For fixed endpoints $(z,z')$ with $z_{1:3}\neq0$, the set $\{g:gz=z'\}$ is a coset of the circle
stabilizing $z$. The states have a density and $\exp$ pushes the Gaussian coefficients to a law with a density on
$SO(3)$, so the conditional law of $g$ on that circle has a density for almost every $(z,z')$. Hence
$\mathbb E\|\hat g-g\|^2\ge\mathbb E\,\operatorname{tr}\operatorname{Cov}(g\mid z,z')>0$ for every endpoint
estimator $\hat g$. $\square$

\paragraph{Remark (what Prop.~\ref{prop:endpoint}b does not show).}
\begin{itemize}
\item \emph{Continuity is essential.} $S$ minus one point is contractible, so $SO(3)\to S^2$ has a continuous
section $s$ there. Ignoring the coefficient bound, $R(z,z')=s(z')s(z)^{-1}$ is a measurable rule that is exact on
every triple avoiding that point, i.e.\ off a null set.
\item \emph{A gap for any fixed Lipschitz class.} Let $J_\epsilon(q)$ be the expected transport and composition loss
on data triples with pairwise angles $<\epsilon$, small enough that the chiral rule is within the bound
($J_\epsilon=0$ for it). For every $L$, the infimum of $J_\epsilon$ over $L$-Lipschitz rules with bounded
coefficients is positive for the true algebra. Otherwise, by Arzel\`a--Ascoli a subsequence converges locally
uniformly to such a rule with $J_\epsilon=0$, by dominated convergence (bounded integrand). The integrand is
continuous and the support contains every such triple on every orbit sphere, so the rule is exact there, which
contradicts (b).
\item \emph{No quantitative statement.} We do not bound this gap, and it may vanish as $L\to\infty$. Nor do we
compare the full training objective, which includes large steps and clipping, between the two algebras.
\end{itemize}

\paragraph{Corollary~\ref{cor:agree}.}
This is Prop.~\ref{prop:endpoint}(a) applied to both models, together with Prop.~\ref{prop:classify}(b) and the
distance $\frac12$ between the diagonal and a chiral span. $\square$
